\hypertarget{hac}{}
\hypertarget{full-cross-portfolio-newey-west-covariance}{%
\chapter{Full cross-portfolio Newey-West
covariance}\label{full-cross-portfolio-newey-west-covariance}}

This chapter builds the object on which all joint inference rests: the
full $25\times25$ covariance of the estimated alpha vector. The general
long-run-variance and Newey--West theory was derived in \S\ref{prelim-hac};
here we motivate why an ordinary --- or even a portfolio-by-portfolio ---
covariance is inadequate for this problem, assemble the joint matrix from
the influence scores of Chapter~\ref{factor-regression-and-the-meaning-of-alpha},
and audit it numerically.

\hypertarget{why-ordinary-ols-covariance-is-inadequate}{%
\section{Why ordinary OLS covariance is
inadequate}\label{why-ordinary-ols-covariance-is-inadequate}}

The classical OLS variance $\sigma^2(X^{\top}X)^{-1}$ rests on residuals
that are homoskedastic, serially independent, and (across portfolios)
mutually uncorrelated. Monthly factor-model residuals violate all three,
so the classical variance is the wrong uncertainty object. The general
form is the sandwich of \S\ref{prelim-ols},
$\operatorname{Var}(\hat\theta)=(X^{\top}X)^{-1}X^{\top}\Omega X(X^{\top}X)^{-1}$,
in which the ``bread'' $(X^{\top}X)^{-1}$ is fixed but the ``meat''
$X^{\top}\Omega X$ must be estimated without assuming $\Omega=\sigma^2 I$.

\hypertarget{heteroskedasticity}{%
\subsection{Heteroskedasticity}\label{heteroskedasticity}}

The conditional residual variance $\operatorname{Var}(\varepsilon_{it}\mid X)$
changes over time --- volatility clusters, with calm and turbulent regimes.
If the true $\operatorname{Var}(\varepsilon_{it})=\sigma_t^2$ varies but the
classical formula plugs in a single $\hat\sigma^2$, the reported standard
error is biased, and the direction of the bias depends on whether large
residuals coincide with high-leverage months. The ARCH-LM test rejects
constant conditional variance for all 25 portfolios (Chapter~13).

\hypertarget{autocorrelation}{%
\subsection{Autocorrelation}\label{autocorrelation}}

Residual shocks remain correlated across nearby months, so
$\operatorname{Cov}(\varepsilon_{it},\varepsilon_{i,t-\ell})\neq 0$ for
small $\ell$. The classical variance uses only the contemporaneous term and
ignores these lag covariances; the correct object is the \emph{long-run}
variance \eqref{eq:lrv}, which sums the autocovariances across lags. Ignoring
positive serial dependence understates the true sampling variability. The
Ljung--Box test rejects the no-autocorrelation null for many portfolios
(Chapter~13), and the Newey--West estimator includes lags to $L=12$.

\hypertarget{cross-sectional-dependence}{%
\subsection{Cross-sectional
dependence}\label{cross-sectional-dependence}}

Different portfolios are hit by the same monthly shocks --- a common market
move affects all 25 at once --- so
$\operatorname{Cov}(\varepsilon_{it},\varepsilon_{jt})\neq 0$ across $i\neq j$.
This is dependence \emph{across the cross-section within a date}, invisible
to any single-portfolio calculation. It is precisely the off-diagonal
information that the joint test of Chapter~7 and the multivariate shrinkage
of Chapter~8 consume: the variance of a difference
$\hat\alpha_i-\hat\alpha_j$ is $V_{ii}+V_{jj}-2V_{ij}$, and dropping $V_{ij}$
misstates every pairwise comparison.

Twenty-five separate HAC standard errors preserve heteroskedasticity and
serial dependence \emph{within} each portfolio --- they give the 25 diagonal
entries of the covariance --- but discard the cross-portfolio covariances.
The construction below recovers the full matrix.

\hypertarget{influence-score-construction}{%
\section{Influence-score
construction}\label{influence-score-construction}}

Because every portfolio uses the same factor matrix $X$, the OLS intercept
estimation error is a weighted sum of residuals with weights common to all
portfolios (\S\ref{deriving-the-intercept-influence-weights}). Let $a_t$ be
the intercept weight for month $t$, $e_t$ the 25-vector of same-month
residuals, and $g_t=a_t e_t$ the 25-vector alpha score. The lag-$\ell$
score cross-product and the Bartlett weights are
\begin{equation}\label{eq:ch6-gamma}
\Gamma_\ell = \sum_{t=\ell+1}^{T} g_t\,g_{t-\ell}^{\top},
\qquad w_\ell = 1 - \frac{\ell}{L+1},
\end{equation}
and the joint HAC covariance is
\begin{equation}\label{eq:ch6-vhac}
\hat V_{\mathrm{HAC}} = c\,\Bigl\{\Gamma_0
+ \sum_{\ell=1}^{L} w_\ell\,(\Gamma_\ell + \Gamma_\ell^{\top})\Bigr\},
\end{equation}
with $L=12$ lags and the finite-sample correction $c=T/(T-P)$, $P=K+1$.
This is the applied form of the general estimator \eqref{eq:nw} and its
joint specialisation \eqref{eq:jointhac} derived in \S\ref{prelim-hac}; the
Bartlett weights decline linearly with lag and guarantee that
$\hat V_{\mathrm{HAC}}$ is positive semi-definite (the Fej\'er-kernel
argument there).

\hypertarget{deriving-the-joint-covariance-from-first-principles}{%
\section{Deriving the joint covariance from first
principles}\label{deriving-the-joint-covariance-from-first-principles}}

The intercept influence weight is itself
$a_t = q_0^{\top}(X^{\top}X)^{-1}x_t$, where $q_0=(1,0,\dots,0)^{\top}$
selects the intercept and $x_t^{\top}$ is row $t$ of $X$
(\S\ref{deriving-the-intercept-influence-weights}); it is a fixed scalar
depending only on the common design, not on any portfolio's returns. Stack
the $N$ portfolio residuals at month $t$ into
$e_t=(\varepsilon_{1t},\dots,\varepsilon_{Nt})^{\top}$. Because the same
$a_t$ multiplies every portfolio, the $N$-vector of alpha-estimation errors
and its per-month score are
\begin{equation}\label{eq:ch6-alphaerr}
\hat\alpha-\alpha = \sum_{t=1}^{T} a_t e_t = \sum_{t=1}^{T} g_t,
\qquad g_t = a_t e_t .
\end{equation}
Without assuming independence through time, the exact covariance is the
full double sum
\begin{equation}\label{eq:ch6-doublesum}
\operatorname{Var}(\hat\alpha-\alpha)
= \sum_{t=1}^{T}\sum_{s=1}^{T}\operatorname{Cov}(g_t,g_s).
\end{equation}

\textbf{Grouping by lag.} The double sum has $T^2$ terms indexed by
$(t,s)$; re-index them by the lag $\ell=t-s$. Under weak stationarity the
covariance depends on $\ell$ alone, so define
$\Gamma_\ell=\operatorname{Cov}(g_t,g_{t-\ell})=\mathbb{E}[g_t g_{t-\ell}^{\top}]$.
The $T$ diagonal terms ($\ell=0$) contribute $\Gamma_0$. For each
$\ell\ge1$, the pairs with $t-s=\ell$ contribute $\Gamma_\ell$, while the
pairs with $s-t=\ell$ contribute
$\operatorname{Cov}(g_t,g_{t+\ell})=\Gamma_{-\ell}=\Gamma_\ell^{\top}$
(covariance of a stationary series satisfies $\Gamma_{-\ell}=\Gamma_\ell^{\top}$).
Collecting the two directions, the double sum collapses to the one-sided
lag representation
\begin{equation}\label{eq:ch6-lagrep}
\operatorname{Var}(\hat\alpha-\alpha)
= \Gamma_0 + \sum_{\ell=1}^{T-1}\bigl(\Gamma_\ell+\Gamma_\ell^{\top}\bigr),
\end{equation}
which is exactly the long-run covariance \eqref{eq:lrv}. The transpose term
is not decoration: it supplies the negative-lag contributions and makes the
matrix symmetric.

\textbf{From population to estimator.} Three finite-sample steps turn the
population identity \eqref{eq:ch6-lagrep} into the usable estimator
\eqref{eq:ch6-vhac}. First, each population $\Gamma_\ell$ is replaced by its
sample cross-product $\hat\Gamma_\ell=\sum_{t=\ell+1}^{T} g_t g_{t-\ell}^{\top}$
from \eqref{eq:ch6-gamma}. Second, the sum is \emph{truncated} at a
bandwidth $L\ll T$, because $\hat\Gamma_\ell$ at large $\ell$ is built from
only $T-\ell$ products and is too noisy to include. Third, each retained lag
is \emph{tapered} by the Bartlett weight $w_\ell=1-\ell/(L+1)$. Element
$(i,j)$ of the result estimates
$\operatorname{Cov}(\hat\alpha_i,\hat\alpha_j)$, and setting $i=j$ recovers
the single-portfolio HAC variance \eqref{eq:hac-alpha} --- the identity the
numerical audit below exploits.

\begin{longtable}[]{@{}p{0.20\textwidth}p{0.34\textwidth}p{0.36\textwidth}@{}}
\toprule\noalign{}
Object & Dimension & Role \\
\midrule\noalign{}
\endhead
\bottomrule\noalign{}
\endlastfoot
X & T \ensuremath{\times } (K+1) & Common regression design \\
a & T \ensuremath{\times } 1 & Intercept influence weights \\
e\textsubscript{t} & N \ensuremath{\times } 1 & Same-month residual vector \\
g\textsubscript{t}=a\textsubscript{t}e\textsubscript{t} & N \ensuremath{\times } 1 &
Same-month alpha score \\
\ensuremath{\Gamma }\textsubscript{\ensuremath{\ell }} & N \ensuremath{\times } N & Lag-\ensuremath{\ell } score cross-product \\
\ensuremath{\hat{V}}\textsubscript{HAC} & N \ensuremath{\times } N & Joint alpha-estimation covariance \\
\end{longtable}

\section{Tapering, lag length, and the finite-sample correction}\label{ch6-bandwidth}

Three choices in \eqref{eq:ch6-vhac} deserve explanation, because they
control whether the estimate is a valid covariance and how much it can be
trusted.

\textbf{Why taper (positive semi-definiteness).} A truncated \emph{un}weighted
sum $\Gamma_0+\sum_{\ell=1}^{L}(\Gamma_\ell+\Gamma_\ell^{\top})$ --- a
rectangular window --- need not be positive semi-definite, and a
``covariance'' with a negative eigenvalue is meaningless (it would imply a
portfolio combination of alphas with negative variance). The Bartlett taper
fixes this. The weighted sum can be written as a
$T^{-1}\sum_{|\ell|\le L}(1-|\ell|/(L+1))\sum_t g_t g_{t-\ell}^{\top}$, which
is a \emph{Fej\'er-kernel} (triangular) smoothing of the periodogram; the
Fej\'er kernel is non-negative, so the resulting matrix is guaranteed
positive semi-definite for any data (\S\ref{prelim-hac}). This is why the
linearly declining weights $w_\ell=1-\ell/(L+1)$ are used rather than a
sharp cutoff.

\textbf{Choosing the bandwidth $L$.} The lag length trades bias against
variance. Too small an $L$ omits genuine dependence beyond lag $L$ and
biases the variance downward when autocorrelation is positive and
persistent; too large an $L$ admits many noisy $\hat\Gamma_\ell$ and
inflates the estimator's own sampling variability. Consistency of the HAC
estimator requires $L\to\infty$ but $L/T\to0$ as $T\to\infty$, so the
bandwidth must grow with the sample yet slowly. The canonical monthly
regressions fix $L=12$: one year of lags, a standard convention that lets
residual dependence persist up to twelve months while keeping each
$\hat\Gamma_\ell$ well estimated on $T=1{,}193$ months. Where the series is
short --- the annual forward-validation statistics of Chapter~11 --- a fixed
$L=12$ would be too large, so the automatic Newey--West (1994) rule
$L=\lfloor 4(T/100)^{2/9}\rfloor$ is used instead (\S\ref{prelim-hac}).

\textbf{The finite-sample correction $c=T/(T-P)$.} Residuals are computed
from an estimated $\hat\theta$, which consumes $P=K+1$ degrees of freedom
and makes the raw residual cross-products slightly too small on average,
exactly as dividing by $T$ rather than $T-P$ understates the OLS residual
variance. Multiplying by $c=T/(T-P)$ is the multivariate analogue of that
degrees-of-freedom correction and removes the leading small-sample downward
bias; with $T=1{,}193$ and $P=4$ it is $c\approx1.003$, negligible here but
retained for exactness and reported in the manifest.

\hypertarget{toy-covariance-calculation}{%
\section{Toy covariance calculation}\label{toy-covariance-calculation}}

Reusing the five-month example of Chapter~\ref{factor-regression-and-the-meaning-of-alpha},
the centred factor gives $a_t=0.2$ for every $t$. Ignoring lag terms only
for illustration, the residual cross-product matrix across the three
portfolios A, B, C is
\begin{equation}\label{eq:ch6-toy}
E^{\top}E = \begin{bmatrix}
 0.140 &  0.130 & -0.022\\
 0.130 &  0.123 & -0.017\\
-0.022 & -0.017 &  0.0086
\end{bmatrix}.
\end{equation}
Thus the contemporaneous score cross-product is $\Gamma_0 = 0.2^2\,E^{\top}E$,
and with $P=2$ regression parameters the finite-sample factor is
$c=T/(T-P)=5/3$, so the illustrative covariance is
$\tfrac53\times 0.04\times E^{\top}E$. The positive A--B entry means the two
portfolios' alpha-estimation errors tend to move together; treating them as
independent would discard exactly this cross-portfolio information.

\hypertarget{numerical-audit}{%
\section{Numerical audit}\label{numerical-audit}}

A HAC covariance is only useful if it is genuinely a covariance. Bartlett
weights guarantee positive semi-definiteness in theory
(\S\ref{prelim-hac}), but finite-precision arithmetic can introduce tiny
asymmetries and negative eigenvalues, so the estimate is checked before
use: it is symmetrised, its eigenvalues are inspected, and only negligible,
scale-relative negatives are floored. A materially indefinite result
\emph{raises} rather than being silently regularised, so numerical repair
is always visible in the run manifest.

\begin{verbatim}
V = 0.5 * (V + V.T)                   # remove floating asymmetry
eig = np.linalg.eigvalsh(V)
scale = max(np.linalg.norm(V, 2), 1.0)
tol = 1e-12 * scale
if eig.min() < -tol:
    raise ArithmeticError("Materially indefinite HAC covariance")
condition_number = np.linalg.cond(V)
# Record any tiny numerical flooring; never hide large regularisation.
\end{verbatim}

\textbf{Paper result.} The canonical $25\times25$ matrix is positive
definite without any flooring and has condition number $139.4$; the ratio
of largest to smallest eigenvalue is finite and manageable, so the linear
solves behind the Wald statistic (Chapter~7) and the empirical-Bayes
posterior (Chapter~8) are numerically stable. A large condition number
would signal near-collinearity of the alpha errors and unstable inversion;
$139.4$ is comfortably below that regime.

\textbf{Independent replication.} The diagonal of $\hat V_{\mathrm{HAC}}$
must equal the 25 single-portfolio HAC variances \eqref{eq:hac-alpha}.
Recomputing each portfolio's alpha and HAC standard error with an
independent implementation (\texttt{statsmodels}) reproduces the diagonal
to a tolerance of $10^{-12}$ for all 25 portfolios. This checks the joint
construction against a separate codebase; two functions calling the same
internal routine would not constitute independent replication. Off-diagonal
entries have no single-portfolio analogue and are validated instead through
the positive-definiteness and symmetry checks above.

\hypertarget{what-hac-does-and-does-not-do}{%
\section{What HAC does and does not
do}\label{what-hac-does-and-does-not-do}}

It is worth stating the estimator's remit precisely, because HAC is easy to
over-trust. It makes the \emph{covariance estimate} robust to the forms of
dependence it models --- heteroskedasticity, serial correlation to lag $L$,
and contemporaneous cross-portfolio correlation --- and thereby makes the
joint Wald test of Chapter~7 asymptotically valid without an IID or
Gaussian assumption. It does none of the following: it does not make the
factor model correct, does not recover omitted priced factors or repair
time-varying betas, does not capture serial dependence beyond lag $L$, and
does not fix data errors such as a bad join or the wrong units. HAC is a
statement about the \emph{sampling distribution of the estimator}, not
about the \emph{truth of the model}; a portfolio can have a perfectly
estimated HAC standard error around an alpha that is entirely an artefact
of an incomplete benchmark.

\begin{longtable}[]{@{}p{0.29\textwidth}p{0.63\textwidth}@{}}
\toprule\noalign{}
HAC helps with & HAC does not repair \\
\midrule\noalign{}
\endhead
\bottomrule\noalign{}
\endlastfoot
Heteroskedasticity and finite-lag serial dependence in covariance
estimation & Omitted factors, nonlinear payoffs or time-varying betas \\
Cross-portfolio covariance when scores are stacked by common date & Bad
joins, wrong units or an incorrect data vintage \\
Asymptotically robust Wald inference & Gaussian likelihood assumptions
used elsewhere \\
\end{longtable}

